\clearpage\section{MOS transistor schematics}
Source: SCH in \texttt{01\_hysteresis.sp}.\par
The ALT comparison is from \texttt{05\_alternative\_dc.sp}. Both topologies are shown because their feedback connections differ. Here vp is $V_{DD}$ and vn is ground.

The drawings use the device and node names of the cited netlists. Equal net labels denote an electrical connection. Unless stated otherwise, NMOS bulk terminals connect to ground and PMOS bulk terminals connect to $V_{DD}$. The bulk wires are omitted from the three-terminal symbols for readability.
\begin{figure}[!htbp]\centering
\resizebox{7.3cm}{!}{\begin{circuitikz}[american,font=\small]
\node[pmos] (p1) at (0,4) {MP1};\node[pmos] (p2) at (0,2) {MP2};
\node[nmos] (n2) at (0,-1) {MN2};\node[nmos] (n1) at (0,-3) {MN1};
\draw (p1.S) -- (0,5) node[above] {vp};\draw (p1.D) -- (p2.S) node[midway,left] {p};
\draw (p2.D) -- (0,.5) node[circ] {} -- (n2.D);\draw (0,.5) -- (2,.5) node[right] {y};
\draw (n2.S) -- (n1.D) node[midway,left] {n};\draw (n1.S) -- (0,-4) node[below] {vn};
\foreach \t in {p1,p2,n2,n1}{\draw (\t.G) -- ++(-.8,0) node[left] {a};}
\node[pmos] (p3) at (4,2) {MP3};\node[nmos] (n3) at (4,-1) {MN3};
\draw (p3.S) -- (4,3.2) -- (0,3.2) node[circ] {};
\draw (n3.S) -- (4,-2.2) -- (0,-2.2) node[circ] {};
\draw (p3.D) -- ++(1.1,0) node[right] {vn};\draw (n3.D) -- ++(1.1,0) node[right] {vp};
\draw (p3.G) -- ++(-.7,0) node[left] {y};\draw (n3.G) -- ++(-.7,0) node[left] {y};
\end{circuitikz}}
\caption{Complete six-MOS SCH. MP1/MP2: $W=3.6\,\mu$m; MN1/MN2: $W=1.2\,\mu$m; MP3: $W=1.8f\,\mu$m; MN3: $W=0.6f\,\mu$m. All $L=0.25\,\mu$m. Feedback source/drain assignments retain the netlist: MP3 $(D,G,S,B)=(vn,y,p,vp)$; MN3 $(D,G,S,B)=(vp,y,n,vn)$.}
\end{figure}
\begin{figure}[!htbp]\centering
\resizebox{14.5cm}{!}{\begin{circuitikz}[american,font=\small]

\begin{scope}[shift={(0,0)}]
\node[pmos] (XINp) at (0,1) {};
\node[nmos] (XINn) at (0,-1) {};
\draw (XINp.S) -- (0,2) node[above] {$V_{DD}$};
\draw (XINn.S) -- (0,-2) node[ground] {};
\draw (XINp.D) -- (0,0) -- (XINn.D);
\draw (XINp.G) -- (-1.15,1) -- (-1.15,-1) -- (XINn.G);
\draw (-1.15,0) node[circ] {} -- (-1.8,0) node[above] {a};
\draw (0,0) node[circ] {} -- (0.85,0) node[above] {y};
\node[right] at (0.25,1) {MP};
\node[right] at (0.25,-1) {MN};
\node at (0,2.6) {XIN};
\coordinate (XINin) at (-1.8,0);
\coordinate (XINout) at (0.85,0);
\end{scope}

\begin{scope}[shift={(5,0)}]
\node[pmos] (XSENSEp) at (0,1) {};
\node[nmos] (XSENSEn) at (0,-1) {};
\draw (XSENSEp.S) -- (0,2) node[above] {$V_{DD}$};
\draw (XSENSEn.S) -- (0,-2) node[ground] {};
\draw (XSENSEp.D) -- (0,0) -- (XSENSEn.D);
\draw (XSENSEp.G) -- (-1.15,1) -- (-1.15,-1) -- (XSENSEn.G);
\draw (-1.15,0) node[circ] {} -- (-1.8,0) node[above] {y};
\draw (0,0) node[circ] {} -- (0.85,0) node[above] {q};
\node[right] at (0.25,1) {MP};
\node[right] at (0.25,-1) {MN};
\node at (0,2.6) {XSENSE};
\coordinate (XSENSEin) at (-1.8,0);
\coordinate (XSENSEout) at (0.85,0);
\end{scope}

\begin{scope}[shift={(10,0)}]
\node[pmos] (XFBp) at (0,1) {};
\node[nmos] (XFBn) at (0,-1) {};
\draw (XFBp.S) -- (0,2) node[above] {$V_{DD}$};
\draw (XFBn.S) -- (0,-2) node[ground] {};
\draw (XFBp.D) -- (0,0) -- (XFBn.D);
\draw (XFBp.G) -- (-1.15,1) -- (-1.15,-1) -- (XFBn.G);
\draw (-1.15,0) node[circ] {} -- (-1.8,0) node[above] {q};
\draw (0,0) node[circ] {} -- (0.85,0) node[above] {y};
\node[right] at (0.25,1) {MP};
\node[right] at (0.25,-1) {MN};
\node at (0,2.6) {XFB};
\coordinate (XFBin) at (-1.8,0);
\coordinate (XFBout) at (0.85,0);
\end{scope}
\draw (XINout) -- (XSENSEin);\draw (XSENSEout) -- (XFBin);
\draw (XFBout) -- (11.5,0) -- (11.5,-3) -- (1.8,-3) -- (1.8,0) node[circ] {};
\end{circuitikz}}
\caption{ALT: strong input inverter XIN ($W_n/W_p=1.2/3.6\,\mu$m) and two weaker cross-coupled inverters XSENSE/XFB ($0.6/1.8\,\mu$m). The XFB output connects back to y; all lengths are $0.25\,\mu$m.}
\end{figure}
